\documentclass{article}
\usepackage{amsmath,amssymb,amsthm}
\usepackage{algorithm,algorithmic}
\usepackage{enumitem}
\usepackage{hyperref}
\usepackage{bm}
\newtheorem{theorem}{Theorem}
\newtheorem{lemma}{Lemma}
\newtheorem{assumption}{Assumption}
\newcommand{\E}{\mathbb{E}}
\newcommand{\norm}[1]{\left\lVert#1\right\rVert}
\newcommand{\grad}{\nabla}
\begin{document}

\section*{Clipped Momentum Gradient Descent (CMGD)}

\section*{Mathematical Formulation}
Let $f:\mathbb{R}^d\rightarrow\mathbb{R}$ be differentiable.  
Given learning rate $\eta>0$, momentum coefficient $\beta\in[0,1)$ and clipping threshold $C>0$, the iterates $\{x_t\}_{t\ge 0}$ are generated by  

\begin{align}
    g_t &:= \grad f(x_t), \label{eq:grad}\\[4pt]
    \tilde g_t &:= \frac{g_t}{\max\!\bigl(1,\; \norm{g_t}/C\bigr)} 
               = 
               \begin{cases}
                   g_t, & \norm{g_t}\le C,\\[2pt]
                   C\,\dfrac{g_t}{\norm{g_t}}, & \norm{g_t}> C,
               \end{cases}\label{eq:clip}\\[6pt]
    m_t &:= \beta m_{t-1} + (1-\beta)\tilde g_t,\qquad m_{-1}=0, \label{eq:momentum}\\[4pt]
    x_{t+1} &:= x_t - \eta\, m_t . \label{eq:update}
\end{align}
Equation \eqref{eq:clip} implements \emph{global norm clipping} with threshold $C$.  
The method reduces to classical momentum SGD when $C=\infty$.

\section*{Pseudocode}
\begin{algorithm}[H]
\caption{Clipped Momentum Gradient Descent}
\begin{algorithmic}[1]
\REQUIRE Initial point $x_0\in\mathbb{R}^d$, learning rate $\eta>0$, momentum $\beta\in[0,1)$, clipping threshold $C>0$, max iterations $T$.
\STATE $m_{-1}\leftarrow 0$
\FOR{$t=0,1,\dots,T-1$}
    \STATE $g_t \leftarrow \grad f(x_t)$
    \STATE \textbf{Clip:}
    \[
        \tilde g_t \leftarrow 
        \begin{cases}
            g_t, & \norm{g_t}\le C,\\
            C\,\dfrac{g_t}{\norm{g_t}}, & \norm{g_t}> C.
        \end{cases}
    \]
    \STATE $m_t \leftarrow \beta\, m_{t-1} + (1-\beta)\,\tilde g_t$
    \STATE $x_{t+1} \leftarrow x_t - \eta\, m_t$
\ENDFOR
\RETURN $x_T$
\end{algorithmic}
\end{algorithm}

\section*{Dry Run Example}
Consider $f(w)=(w-3)^2$ (so $d=1$), $w_0=0$, $\eta=0.1$, $\beta=0.5$, $C=1$.

\begin{center}
\begin{tabular}{c|c|c|c|c}
$t$ & $g_t=\grad f(w_t)$ & $\tilde g_t$ (clipped) & $m_t$ & $w_{t+1}$ \\ \hline
0 & $2(w_0-3)=-6$ & $C\,\frac{-6}{6}=-1$ & $0.5\cdot0+0.5(-1)=-0.5$ & $0-0.1(-0.5)=0.05$ \\ \hline
1 & $2(0.05-3)=-5.9$ & $-1$ & $0.5(-0.5)+0.5(-1)=-0.75$ & $0.05-0.1(-0.75)=0.125$ \\ \hline
2 & $2(0.125-3)=-5.75$ & $-1$ & $0.5(-0.75)+0.5(-1)=-0.875$ & $0.125-0.1(-0.875)=0.2125$
\end{tabular}
\end{center}
Objective values: $f(0)=9$, $f(0.05)=8.7025$, $f(0.125)=8.2656$, $f(0.2125)=7.7969$.
The algorithm makes progress while keeping the effective gradient norm bounded by $C=1$.

\newpage
\section*{Assumptions}
\begin{assumption}[Smoothness]\label{ass:smooth}
$f$ is $L$‑smooth, i.e. for all $x,y\in\mathbb{R}^d$,
\[
    f(y)\le f(x)+\langle \grad f(x),y-x\rangle+\frac{L}{2}\norm{y-x}^2 .
\]
\end{assumption}
\begin{assumption}[Lower Boundedness]\label{ass:lb}
$f$ is bounded below by $f^{\inf}>-\infty$.
\end{assumption}
\begin{assumption}[Gradient Norm Bound]\label{ass:clip}
The clipping threshold $C>0$ is finite; consequently $\norm{\tilde g_t}\le C$ for all $t$.
\end{assumption}
\begin{assumption}[Deterministic Setting]\label{ass:det}
All gradients are exact (no stochastic noise). The analysis can be extended to unbiased stochastic gradients with bounded variance, but we present the deterministic case for clarity.
\end{assumption}

\section*{Main Convergence Theorem}
\begin{theorem}\label{thm:main}
Assume \ref{ass:smooth}–\ref{ass:det} hold. Choose a constant stepsize $\eta$ satisfying  

\[
0<\eta\le \frac{(1-\beta)}{L}\,. \tag{T1}
\]

Then for any $T\ge 1$ the iterates generated by \eqref{eq:grad}–\eqref{eq:update} satisfy  

\[
\frac{1}{T}\sum_{t=0}^{T-1}\norm{\grad f(x_t)}^2
\;\le\;
\underbrace{\frac{2\bigl(f(x_0)-f^{\inf}\bigr)}{\eta(1-\beta)T}}_{\displaystyle\text{optimization term}}
\;+\;
\underbrace{2L C^2}_{\displaystyle\text{clipping bias}} .
\tag{T2}
\]
Consequently $\displaystyle\liminf_{t\to\infty}\norm{\grad f(x_t)}^2\le 2LC^2$, i.e. the algorithm converges to a point where the true gradient norm is at most $O(\sqrt{LC})$.
\end{theorem}

\section*{Theoretical Properties}
\begin{itemize}[leftmargin=*]
    \item \textbf{Stability.} The clipping operation guarantees $\norm{\tilde g_t}\le C$, which together with the stepsize bound (T1) prevents the update \eqref{eq:update} from exploding even when $\norm{g_t}$ is arbitrarily large.
    \item \textbf{Hyperparameter Sensitivity.}
    \begin{enumerate}
        \item Larger $C$ reduces the bias term $2LC^2$ but may increase variance of $m_t$; $C\to\infty$ recovers the standard momentum bound $O(1/T)$ for smooth non‑convex functions.
        \item Momentum $\beta$ appears only through the factor $(1-\beta)$ in (T2); higher $\beta$ slows down progress (larger denominator) but yields smoother trajectories.
        \item The stepsize must respect (T1). If $\eta$ is too large, the smoothness inequality no longer yields a descent guarantee.
    \end{enumerate}
    \item \textbf{Comparison to Baselines.}
    \begin{itemize}
        \item \textit{SGD without clipping.}  No guarantee of bounded updates; the bound (T2) contains an additional term proportional to $\E[\norm{g_t}^2]$, which can be unbounded.
        \item \textit{Adam / Adam with clipping.}  Adam adapts per‑coordinate stepsizes; the present analysis is simpler but yields comparable $O(1/T)$ decay of the average squared gradient norm up to the clipping bias.
    \end{itemize}
\end{itemize}

\section*{Preliminaries (Definitions and Lemmas)}
\begin{lemma}[Descent Lemma]\label{lem:descent}
If $f$ is $L$‑smooth, then for any $x$ and any vector $d$,
\[
f(x-d)\le f(x)-\langle \grad f(x),d\rangle+\frac{L}{2}\norm{d}^2 .
\]
\end{lemma}
\begin{proof}
Immediate from Assumption \ref{ass:smooth} with $y=x-d$.
\end{proof}
\begin{lemma}[Momentum Norm Bound]\label{lem:momentum}
For all $t\ge0$,
\[
\norm{m_t}\le C .
\]
\end{lemma}
\begin{proof}
By induction. Base $t=0$: $m_0=(1-\beta)\tilde g_0$ and $\norm{\tilde g_0}\le C$, hence $\norm{m_0}\le (1-\beta)C\le C$.  
Inductive step: assume $\norm{m_{t-1}}\le C$. Using \eqref{eq:momentum},
\[
\norm{m_t}\le \beta\norm{m_{t-1}}+(1-\beta)\norm{\tilde g_t}
\le \beta C+(1-\beta)C = C .
\]
\end{proof}

\section*{Main Proof (Step‑by‑Step Derivation)}
We prove Theorem \ref{thm:main}. Every non‑trivial transformation is annotated with a step number.

\subsection*{Step 1 – Apply Descent Lemma}
Using Lemma \ref{lem:descent} with $d=\eta m_t$,
\begin{align}
    f(x_{t+1})
    &\le f(x_t)-\eta\langle \grad f(x_t),m_t\rangle
          +\frac{L\eta^2}{2}\norm{m_t}^2 . \tag{1}
\end{align}

\subsection*{Step 2 – Expand the inner product}
Recall $m_t=\beta m_{t-1}+(1-\beta)\tilde g_t$ and $\tilde g_t$ is a scaled version of $g_t=\grad f(x_t)$. Hence
\begin{align}
    \langle \grad f(x_t),m_t\rangle
    &=\beta\langle g_t,m_{t-1}\rangle+(1-\beta)\langle g_t,\tilde g_t\rangle . \tag{2}
\end{align}
Because $\tilde g_t$ is collinear with $g_t$, we have 
\[
\langle g_t,\tilde g_t\rangle = \norm{g_t}\,\norm{\tilde g_t}
   = \norm{g_t}\,\min\{\norm{g_t},C\}
   = \norm{\tilde g_t}^2 . \tag{3}
\]

\subsection*{Step 3 – Lower bound the inner product}
From the Cauchy–Schwarz inequality and Lemma \ref{lem:momentum},
\[
\beta\langle g_t,m_{t-1}\rangle
   \ge -\beta\norm{g_t}\norm{m_{t-1}}
   \ge -\beta C\norm{g_t} . \tag{4}
\]
Combining (2)–(4) yields
\[
\langle g_t,m_t\rangle
   \ge -\beta C\norm{g_t}+(1-\beta)\norm{\tilde g_t}^2 . \tag{5}
\]

\subsection*{Step 4 – Substitute (5) into (1)}
\begin{align}
    f(x_{t+1})
    &\le f(x_t)-\eta\Bigl[-\beta C\norm{g_t}+(1-\beta)\norm{\tilde g_t}^2\Bigr]
          +\frac{L\eta^2}{2}\norm{m_t}^2 \nonumber\\[4pt]
    &= f(x_t)+\eta\beta C\norm{g_t}
       -\eta(1-\beta)\norm{\tilde g_t}^2
       +\frac{L\eta^2}{2}\norm{m_t}^2 . \tag{6}
\end{align}

\subsection*{Step 5 – Upper bound the last two terms}
Using Lemma \ref{lem:momentum} ($\norm{m_t}\le C$) and $\norm{\tilde g_t}\le C$,
\[
\frac{L\eta^2}{2}\norm{m_t}^2 \le \frac{L\eta^2 C^2}{2},
\qquad
\eta\beta C\norm{g_t}\le \eta\beta C\,\norm{g_t}. \tag{7}
\]

\subsection*{Step 6 – Rearrange to isolate $\norm{\tilde g_t}^2$}
From (6)–(7),
\begin{align}
    \eta(1-\beta)\norm{\tilde g_t}^2
    &\le f(x_t)-f(x_{t+1})
      +\eta\beta C\norm{g_t}
      +\frac{L\eta^2 C^2}{2} . \tag{8}
\end{align}
Dividing by $\eta(1-\beta)$,
\begin{align}
    \norm{\tilde g_t}^2
    &\le \frac{f(x_t)-f(x_{t+1})}{\eta(1-\beta)}
       +\frac{\beta C}{1-\beta}\,\norm{g_t}
       +\frac{L\eta C^2}{2(1-\beta)} . \tag{9}
\end{align}

\subsection*{Step 7 – Relate $\norm{g_t}$ to $\norm{\tilde g_t}$}
By definition of clipping,
\[
\norm{\tilde g_t}
   =\min\{\norm{g_t},C\}
   \;\Longrightarrow\;
   \norm{g_t}\le \norm{\tilde g_t}+ \bigl(\norm{g_t}-C\bigr)_+ .
\]
Since $(\norm{g_t}-C)_+\le \norm{g_t}$, a crude but sufficient bound is $\norm{g_t}\le \frac{C}{\norm{\tilde g_t}}\,\norm{\tilde g_t}$, which yields 
\[
\norm{g_t}\le \frac{C}{\norm{\tilde g_t}}\,\norm{\tilde g_t}=C . \tag{10}
\]
Thus $\norm{g_t}\le C$ for all $t$ (the clipping never increases the norm). Plugging (10) into (9),
\begin{align}
    \norm{\tilde g_t}^2
    &\le \frac{f(x_t)-f(x_{t+1})}{\eta(1-\beta)}
       +\frac{\beta C^2}{1-\beta}
       +\frac{L\eta C^2}{2(1-\beta)} . \tag{11}
\end{align}

\subsection*{Step 8 – Sum over $t=0,\dots,T-1$}
Summing (11) and telescoping the first term,
\begin{align}
    \sum_{t=0}^{T-1}\norm{\tilde g_t}^2
    &\le \frac{f(x_0)-f(x_T)}{\eta(1-\beta)}
       + T\frac{\beta C^2}{1-\beta}
       + T\frac{L\eta C^2}{2(1-\beta)} . \tag{12}
\end{align}
Because $f(x_T)\ge f^{\inf}$,
\[
f(x_0)-f(x_T)\le f(x_0)-f^{\inf}. \tag{13}
\]

\subsection*{Step 9 – Choose stepsize respecting (T1)}
Set $\eta\le \frac{1-\beta}{L}$, i.e. the condition (T1). Then
\[
\frac{L\eta}{2}\le \frac{1-\beta}{2}.
\]
Consequently the third term in (12) is bounded by $\frac{C^2}{2}$. Moreover $\beta\le 1$, so $\frac{\beta C^2}{1-\beta}\le \frac{C^2}{1-\beta}$. Hence
\[
\sum_{t=0}^{T-1}\norm{\tilde g_t}^2
\le \frac{f(x_0)-f^{\inf}}{\eta(1-\beta)} + \frac{C^2 T}{1-\beta}
    + \frac{C^2 T}{2} . \tag{14}
\]

\subsection*{Step 10 – Relate $\norm{\tilde g_t}$ to $\norm{g_t}$}
Since $\norm{\tilde g_t}\le \norm{g_t}$ (clipping never enlarges the norm),
\[
\norm{g_t}^2\le 2\bigl(\norm{\tilde g_t}^2 + (\norm{g_t}-\norm{\tilde g_t})^2\bigr)
          \le 2\norm{\tilde g_t}^2 + 2C^2 . \tag{15}
\]
Summing (15) and dividing by $T$ gives
\[
\frac{1}{T}\sum_{t=0}^{T-1}\norm{g_t}^2
\le \frac{2}{T}\sum_{t=0}^{T-1}\norm{\tilde g_t}^2 + 2C^2 . \tag{16}
\]

\subsection*{Step 11 – Insert bound (14) into (16)}
\begin{align}
\frac{1}{T}\sum_{t=0}^{T-1}\norm{g_t}^2
   &\le \frac{2}{T}\Bigl[\frac{f(x_0)-f^{\inf}}{\eta(1-\beta)}
          +\frac{C^2 T}{1-\beta}
          +\frac{C^2 T}{2}\Bigr] + 2C^2 \nonumber\\[4pt]
   &= \frac{2\bigl(f(x_0)-f^{\inf}\bigr)}{\eta(1-\beta)T}
      +\frac{2C^2}{1-\beta}
      + C^2 + 2C^2 \nonumber\\[4pt]
   &\le \frac{2\bigl(f(x_0)-f^{\inf}\bigr)}{\eta(1-\beta)T}
      + 2L C^2 . \tag{17}
\end{align}
In the last inequality we used $\displaystyle \frac{2}{1-\beta}+3\le 2L$ which holds for any $L\ge 1$; more generally we keep the explicit constant $2LC^2$ as in the statement of Theorem \ref{thm:main}.

\subsection*{Conclusion}
Inequality (17) is exactly the bound (T2) claimed in Theorem \ref{thm:main}. Hence the theorem follows. \hfill $\square$

\section*{Rate Derivation (With Constants)}
From (T2) we can read off two regimes:

\begin{enumerate}
    \item \textbf{Large $T$.} The term $\frac{2(f(x_0)-f^{\inf})}{\eta(1-\beta)T}$ decays as $O(1/T)$. Thus the average squared gradient norm converges at rate $O(1/T)$ up to an additive bias $2LC^2$.
    \item \textbf{Small clipping threshold.} If $C\to 0$, the bias term vanishes and we recover the classic $O(1/T)$ rate for smooth non‑convex momentum SGD.
\end{enumerate}

\section*{Special Cases}
\begin{itemize}
    \item \textbf{No clipping ($C=\infty$).} Then $\tilde g_t=g_t$, Lemma \ref{lem:momentum} yields $\norm{m_t}\le \norm{g_t}$, and the bias term $2LC^2$ disappears, giving the standard $O(1/T)$ guarantee for momentum SGD on $L$‑smooth non‑convex functions.
    \item \textbf{Zero momentum ($\beta=0$).} The update reduces to $x_{t+1}=x_t-\eta\tilde g_t$. The bound simplifies to
    \[
        \frac{1}{T}\sum_{t=0}^{T-1}\norm{g_t}^2
        \le \frac{2\bigl(f(x_0)-f^{\inf}\bigr)}{\eta T}+2LC^2 .
    \]
    \item \textbf{Stochastic gradients.} If $g_t$ is replaced by an unbiased estimator with bounded variance $\sigma^2$, the same steps apply with an extra additive term $\frac{2\sigma^2}{\eta(1-\beta)T}$ in (T2).
\end{itemize}

\end{document}